\documentclass[11pt,a4paper]{article}
\usepackage[utf8]{vietnam}
\usepackage{amsmath,amssymb}
\usepackage{tabularx}
\usepackage{tikz}
\usepackage{xcolor}
\usepackage{fancyhdr}
\usepackage[top=1.8cm, bottom=2.0cm, left=1.8cm, right=1.8cm]{geometry}

\setlength{\parskip}{3pt}
\setlength{\parindent}{0pt}

\pagestyle{fancy}
\fancyhf{}
\renewcommand{\headrulewidth}{0.4pt}
\lhead{\footnotesize \textbf{CHÙA HƯƠNG SƠN -- ÔN THI VÀO 10}}
\rhead{\footnotesize \textbf{ĐỀ THI THỬ CUỐI HỌC KỲ I (ĐỀ SỐ 04)}}
\cfoot{\thepage}

\begin{document}

\noindent
\begin{tabular*}{\linewidth}{@{\extracolsep{\fill}}cc@{}}
\textbf{SỞ GIÁO DỤC VÀ ĐÀO TẠO} & \textbf{ĐỀ KIỂM TRA CUỐI HỌC KỲ I} \\
\textbf{CHÙA HƯƠNG SƠN} & \textbf{NĂM HỌC 2026 -- 2027} \\
\textbf{ĐỀ THI THỬ SỐ 04} & \textbf{MÔN: TOÁN -- LỚP 9} \\
\textit{(Đề thi gồm 03 trang)} & \textit{Thời gian làm bài: 90 phút (không kể thời gian phát đề)} \\
\end{tabular*}

\vspace{2mm}
\noindent
\rule{\linewidth}{0.8pt}

\vspace{1mm}
\noindent
Họ và tên thí sinh: \dotfill Số báo danh: \dotfill

\vspace{2mm}
\noindent
\textbf{\large I. TRẮC NGHIỆM KHÁCH QUAN (7,0 điểm)}

\vspace{1mm}
\noindent
\textbf{Phần 1. Trắc nghiệm nhiều lựa chọn (3,0 điểm).} \textit{Từ câu 1 đến câu 12, thí sinh chọn đáp án đúng nhất trong 4 đáp án A, B, C, D rồi ghi vào bài làm.}

\vspace{1.5mm}
\textbf{Câu 1.} Giá trị của biểu thức $4\sqrt{12} - \sqrt{48} + \sqrt{3}$ bằng\\
\begin{tabularx}{\linewidth}{XXXX}
\textbf{A.} $5\sqrt{3}$. & \textbf{B.} $3\sqrt{3}$. & \textbf{C.} $4\sqrt{3}$. & \textbf{D.} $-5\sqrt{3}$.
\end{tabularx}

\vspace{1.5mm}
\textbf{Câu 2.} Cho tam giác $ABC$ vuông tại $A$, có góc nhọn $\widehat{C} = \alpha$. Khẳng định nào sau đây là đúng?\\
\begin{tabularx}{\linewidth}{XXXX}
\textbf{A.} $\cos \alpha = \dfrac{AC}{BC}$. & \textbf{B.} $\sin \alpha = \dfrac{AC}{BC}$. & \textbf{C.} $\tan \alpha = \dfrac{AC}{AB}$. & \textbf{D.} $\cot \alpha = \dfrac{AB}{AC}$.
\end{tabularx}

\vspace{1.5mm}
\textbf{Câu 3.} Bất phương trình nào sau đây là bất phương trình bậc nhất một ẩn?\\
\begin{tabularx}{\linewidth}{XXXX}
\textbf{A.} $-2x + 10 > 0$. & \textbf{B.} $0x - 6 \le 0$. & \textbf{C.} $3x^2 - 1 \ge 0$. & \textbf{D.} $\dfrac{4}{x} - 5 < 0$.
\end{tabularx}

\vspace{1.5mm}
\textbf{Câu 4.} Cho $x < y$. Bất đẳng thức nào sau đây luôn đúng?\\
\begin{tabularx}{\linewidth}{XXXX}
\textbf{A.} $-3x > -3y$. & \textbf{B.} $-3x < -3y$. & \textbf{C.} $2x - 1 > 2y - 1$. & \textbf{D.} $x + 4 > y + 4$.
\end{tabularx}

\vspace{1.5mm}
\textbf{Câu 5.} Căn bậc hai số học của số 169 là\\
\begin{tabularx}{\linewidth}{XXXX}
\textbf{A.} $13$. & \textbf{B.} $13$ và $-13$. & \textbf{C.} $-13$. & \textbf{D.} $169$.
\end{tabularx}

\vspace{1.5mm}
\textbf{Câu 6.} Điều kiện xác định của căn thức $\sqrt{4x - 20}$ là\\
\begin{tabularx}{\linewidth}{XXXX}
\textbf{A.} $x \ge 5$. & \textbf{B.} $x \le 5$. & \textbf{C.} $x > 5$. & \textbf{D.} $x < 5$.
\end{tabularx}

\vspace{1.5mm}
\textbf{Câu 7.} Với mọi số thực $x \le 4$, giá trị của biểu thức $\sqrt{(x-4)^2}$ bằng\\
\begin{tabularx}{\linewidth}{XXXX}
\textbf{A.} $4 - x$. & \textbf{B.} $x - 4$. & \textbf{C.} $-(4-x)$. & \textbf{D.} $(x-4)^2$.
\end{tabularx}

\vspace{1.5mm}
\textbf{Câu 8.} Giá trị của biểu thức $\sqrt[3]{-512}$ bằng\\
\begin{tabularx}{\linewidth}{XXXX}
\textbf{A.} $-8$. & \textbf{B.} $8$. & \textbf{C.} $-64$. & \textbf{D.} $64$.
\end{tabularx}

\vspace{1.5mm}
\textbf{Câu 9.} Độ dài cung tròn $n^\circ$ của đường tròn bán kính $R$ được tính theo công thức\\
\begin{tabularx}{\linewidth}{XXXX}
\textbf{A.} $l = \dfrac{\pi R n}{180}$. & \textbf{B.} $l = \dfrac{\pi R n}{360}$. & \textbf{C.} $l = \dfrac{\pi R^2 n}{180}$. & \textbf{D.} $l = \dfrac{2\pi R n}{180}$.
\end{tabularx}

\vspace{1.5mm}
\textbf{Câu 10.} Cho đường tròn $(O; 7\text{ cm})$ và đường thẳng $m$. Khoảng cách từ $O$ đến $m$ là $9\text{ cm}$. Vị trí tương đối của $m$ và $(O)$ là\\
\begin{tabularx}{\linewidth}{XXXX}
\textbf{A.} Không giao nhau. & \textbf{B.} Cắt nhau. & \textbf{C.} Tiếp xúc nhau. & \textbf{D.} Trùng nhau.
\end{tabularx}

\vspace{1.5mm}
\textbf{Câu 11.} Cho hai đường tròn $(O; 8\text{ cm})$ và $(O'; 3\text{ cm})$. Biết khoảng cách $OO' = 11\text{ cm}$. Vị trí tương đối của $(O)$ và $(O')$ là\\
\begin{tabularx}{\linewidth}{XXXX}
\textbf{A.} Tiếp xúc ngoài. & \textbf{B.} Tiếp xúc trong. & \textbf{C.} Cắt nhau. & \textbf{D.} Không giao nhau.
\end{tabularx}

\vspace{1.5mm}
\textbf{Câu 12.} Trên đường tròn $(O)$, số đo cung nhỏ $MN$ bằng $75^\circ$. Số đo của góc ở tâm $\widehat{MON}$ bằng\\
\begin{tabularx}{\linewidth}{XXXX}
\textbf{A.} $75^\circ$. & \textbf{B.} $150^\circ$. & \textbf{C.} $37{,}5^\circ$. & \textbf{D.} $105^\circ$.
\end{tabularx}

\vspace{3mm}
\noindent
\textbf{Phần 2. Trắc nghiệm Đúng -- Sai (2,0 điểm).} \textit{Trong mỗi ý a), b), c), d) ở câu 13 và câu 14, thí sinh chọn đúng hoặc sai và ghi chữ ``Đ'' hoặc ``S'' vào bài làm.}

\vspace{1.5mm}
\textbf{Câu 13.} Xét bài toán thực tế: ``Trong phòng thí nghiệm, một học sinh cần pha trộn hai loại dung dịch muối nồng độ $10\%$ và nồng độ $30\%$ để thu được $400\text{ gam}$ dung dịch muối có nồng độ $25\%$. Gọi khối lượng dung dịch muối $10\%$ và $30\%$ cần lấy lần lượt là $x$ và $y$ (gam, $0 < x, y < 400$).''\\
Xét tính đúng sai của các khẳng định sau:
\begin{center}
\begin{tabularx}{\linewidth}{|X|c|c|}
\hline
\textbf{Khẳng định} & \textbf{Đúng} & \textbf{Sai} \\
\hline
a) Tổng khối lượng hai dung dịch cần pha là $400\text{ gam}$ nên $x + y = 400$. & & \\
\hline
b) Khối lượng muối nguyên chất trong $400\text{ gam}$ dung dịch $25\%$ là $100\text{ gam}$. & & \\
\hline
c) Phương trình biểu thị khối lượng muối nguyên chất là: $0{,}1x + 0{,}3y = 80$. & & \\
\hline
d) Cần dùng $100\text{ gam}$ dung dịch muối $10\%$ và $300\text{ gam}$ dung dịch muối $30\%$. & & \\
\hline
\end{tabularx}
\end{center}

\vspace{2mm}
\textbf{Câu 14.} Cho nửa đường tròn $(O; R)$ đường kính $AB$. Trên cùng một nửa mặt phẳng bờ $AB$, vẽ hai tiếp tuyến $Ax, By$ với nửa đường tròn. Lấy điểm $M$ bất kì trên nửa đường tròn ($M \ne A, B$), tiếp tuyến tại $M$ cắt $Ax$ và $By$ lần lượt tại $C$ và $D$.\\
Xét tính đúng sai của các khẳng định sau:
\begin{center}
\begin{tabularx}{\linewidth}{|X|c|c|}
\hline
\textbf{Khẳng định} & \textbf{Đúng} & \textbf{Sai} \\
\hline
a) Điểm $M$ thuộc nửa đường tròn nên tam giác $AMB$ vuông tại $M$. & & \\
\hline
b) Đoạn thẳng $CD = AC + BD$. & & \\
\hline
c) Góc $\widehat{COD} = 90^\circ$. & & \\
\hline
d) Tích $AC \cdot BD = 2R^2$. & & \\
\hline
\end{tabularx}
\end{center}

\vspace{3mm}
\noindent
\textbf{Phần 3. Trắc nghiệm trả lời ngắn (2,0 điểm).} \textit{Từ câu 15 đến câu 18, thí sinh chỉ ghi kết quả cuối cùng vào bài làm.}

\vspace{1.5mm}
\textbf{Câu 15.} Nghiệm của bất phương trình $5(x + 1) - 3 > 3(2x - 1) + 4$ là $\dots\dots\dots$

\vspace{1.5mm}
\textbf{Câu 16.} Giá trị của biểu thức $K = \sqrt{(\sqrt{8} - 3)^2} + \sqrt{8}$ bằng $\dots\dots\dots$

\vspace{1.5mm}
\noindent
\begin{minipage}[c]{0.65\linewidth}
\textbf{Câu 17.} Một chiếc thang dài $4\text{ m}$ được đặt dựa vào một bức tường thẳng đứng để sửa chữa mái nhà. Để đảm bảo an toàn lao động, thang cần tạo với mặt đất một góc nghiêng khoảng $65^\circ$ (xem hình vẽ bên). Hỏi đỉnh thang chạm tường ở độ cao bao nhiêu mét so với mặt đất? (làm tròn kết quả đến chữ số thập phân thứ hai của mét).
\end{minipage}\hfill
\begin{minipage}[c]{0.32\linewidth}
\centering
\begin{tikzpicture}[scale=0.7]
  % Bức tường
  \fill[gray!30] (0,0) rectangle (0.5,3.8);
  \draw[thick] (0,0) -- (0,3.8) -- (0.5,3.8) -- (0.5,0);
  % Mặt đất
  \draw[thick] (0,0) -- (3.5,0);
  \fill[pattern=north east lines, pattern color=gray!50] (0,-0.2) rectangle (3.5,0);
  % Thang
  \draw[thick, brown!80!black, line width=1.5pt] (0.5,3.4) -- (2.3,0);
  % Các bậc thang
  \foreach \t in {0.2, 0.4, 0.6, 0.8} {
    \draw[thick, brown!80!black] ($(0.5,3.4)!\t!(2.3,0) + (-0.1, 0.05)$) -- ($(0.5,3.4)!\t!(2.3,0) + (0.1, -0.05)$);
  }
  \node[right] at (1.5, 2.0) {$4\text{ m}$};
  \node[left] at (0, 1.7) {$h = ?$};
  % Góc vuông
  \draw (0.5,0) rectangle (0.8,0.3);
  % Góc nghiêng 65 độ
  \draw (2.3,0) ++(180:0.6) arc (180:115:0.6);
  \node at (1.7, 0.35) {\footnotesize $65^\circ$};
\end{tikzpicture}
\end{minipage}

\vspace{1.5mm}
\textbf{Câu 18.} Một miếng kim loại trang trí có dạng hình quạt tròn với bán kính $R = 25\text{ cm}$ và góc ở tâm tương ứng là $72^\circ$. Tính diện tích bề mặt của miếng kim loại đó (lấy $\pi \approx 3{,}14$, đơn vị: $\text{cm}^2$).

\vspace{4mm}
\noindent
\textbf{\large II. TỰ LUẬN (3,0 điểm).} \textit{Thí sinh trình bày lời giải chi tiết vào tờ giấy thi.}

\vspace{2mm}
\textbf{Câu 19 (1,0 điểm).}
\begin{itemize}
  \item[a)] \textit{(0,5 điểm)} Giải hệ phương trình: $\begin{cases} 5x + 2y = 9 \\ 2x - y = 0 \end{cases}$.
  \item[b)] \textit{(0,5 điểm)} Cho biểu thức $M = \left( \dfrac{1}{\sqrt{x}-1} - \dfrac{1}{\sqrt{x}} \right) : \left( \dfrac{\sqrt{x}+1}{\sqrt{x}-2} - \dfrac{\sqrt{x}+2}{\sqrt{x}-1} \right)$ với $x > 0, x \ne 1, x \ne 4$. Rút gọn biểu thức $M$.
\end{itemize}

\vspace{2mm}
\textbf{Câu 20 (1,5 điểm).} Cho nửa đường tròn $(O; R)$ đường kính $AB$. Trên cùng một nửa mặt phẳng bờ $AB$, vẽ hai tiếp tuyến $Ax, By$. Lấy điểm $M$ thuộc nửa đường tròn ($M \ne A, B$), tiếp tuyến tại $M$ cắt $Ax, By$ lần lượt tại $C$ và $D$. Gọi $K$ là giao điểm của $AD$ và $BC$.
\begin{itemize}
  \item[a)] \textit{(0,75 điểm)} Chứng minh rằng $\widehat{COD} = 90^\circ$ và $CD = AC + BD$.
  \item[b)] \textit{(0,75 điểm)} Chứng minh rằng: $MK \parallel AC$ và $MK \cdot (AC + BD) = AC \cdot BD$.
\end{itemize}

\vspace{2mm}
\textbf{Câu 21 (0,5 điểm).}
\begin{itemize}
  \item[a)] \textit{(0,25 điểm)} Từ một tấm bìa hình vuông cạnh $60\text{ cm}$, người ta cắt bỏ ở bốn góc bốn hình vuông nhỏ bằng nhau có cạnh bằng $x$ (cm) rồi gập các mép lại để tạo thành một chiếc hộp hình chữ nhật không có nắp. Tìm $x$ để thể tích của chiếc hộp là lớn nhất.
  \item[b)] \textit{(0,25 điểm)} Cho hai số thực dương $x, y$ thỏa mãn điều kiện $x + y \le 2$. Tìm giá trị nhỏ nhất của biểu thức:
  \[
  S = \dfrac{1}{x^2 + y^2} + \dfrac{1}{xy}.
  \]
\end{itemize}

\vspace{4mm}
\begin{center}
\textbf{---------- HẾT ----------}\\
\textit{(Cán bộ coi thi không giải thích gì thêm)}
\end{center}

\end{document}
